\chapter{Implementation}
\label{app:impl}

\keybox{what this appendix is for}{%
The thesis compares two published methods and two proposed ones. All four run
from a single checkout, differing only by command-line flags, and this appendix
says exactly what each flag does and where in the code it acts. It is the answer
to ``how do I know the 3DGS arm is really 3DGS?'', and it is a code-level answer.}

\section{Four configurations}

Running each method from its own repository would mean the difference between two
numbers contained the difference between two dataloaders, two metric
implementations and two LPIPS backbones. Every comparison in this thesis is
therefore between configurations of one checkout, and the comparison is the
flags.

\begin{center}\small
\begin{tabular}{llp{0.44\linewidth}}
  \toprule
  configuration & flags & what it is \\
  \midrule
  the Mip-Splatting arm & \texttt{--kernel\_size 0.1} &
    Mip-Splatting as published: 3D smoothing filter on, 2D dilation with the
    opacity compensation. \\
  the 3DGS arm & \texttt{--kernel\_size 0.3 --disable\_3D\_filter
                  --disable\_2D\_mip\_compensation} &
    3DGS's band-limit semantics exactly: no world-space filter, a fixed $0.3$
    screen-space dilation, no opacity compensation. \\
  the anisotropic filter & \texttt{--kernel\_size 0.1 --use\_fisher\_filter} &
    Method I. Replaces the scalar 3D filter with the anisotropic one, floored at
    Mip-Splatting's own value. \\
  the angular mask & post-hoc mask, no training flag &
    Method II. Reads a trained model and the training view directions; masks
    spherical-harmonic degrees. \\
  \bottomrule
\end{tabular}
\end{center}

\section{Two switches for the 3DGS arm}

Mip-Splatting has \emph{two} band-limits and they live in different places.

The first is the 3D smoothing filter: a per-primitive scalar variance added in
world space, computed in Python from the training cameras. Zeroing it is a
one-line change, and with \texttt{--kernel\_size 0.3} --- 3DGS's fixed screen-space
dilation --- the model's world-space semantics are 3DGS's.

The second is the screen-space dilation's \emph{opacity compensation},
$\rho = \sqrt{\det \Sigma' / \det(\Sigma' + kI)}$, which lives in CUDA. Vanilla
3DGS dilates the projected covariance by a fixed $0.3$ and applies no such term.
A Python attribute cannot reach it.

For a long time the 3DGS arm carried only the first switch. That is the defect of
Appendix~\ref{app:defects}: the arm labelled 3DGS was still anti-aliasing with
Mip-Splatting's screen-space term, at exactly the scales where the falsification
gate was scored.

\paragraph{What \texttt{--disable\_2D\_mip\_compensation} touches.} The flag is
threaded from \texttt{arguments/\_\_init\_\_.py} through
\texttt{gaussian\_renderer} into the rasteriser's settings tuple, and read in
both the forward and the backward CUDA kernels:

\begin{itemize}[leftmargin=1.4em,itemsep=2pt]
  \item \texttt{forward.cu}, \texttt{computeCov2D}: $\rho$ is computed only when
        the flag is off; otherwise the coefficient is $1$.
  \item \texttt{backward.cu}, \texttt{computeCov2DCUDA}: the same, and the
        gradient path from opacity through the 2D covariance is removed with it.
\end{itemize}

Because the switch reaches the binary, the pre-flight asserts that the
\emph{compiled} rasteriser carries the field --- checking the source would prove
nothing about what was built.

\section{The degeneracy guard}

Both kernels carry a numerical guard: a primitive whose projected covariance is
singular to the tolerance the determinants are clamped at contributes nothing.
That guard originally sat \emph{inside} the compensation branch, so switching the
compensation off removed the guard as well --- which is not what the switch was
supposed to isolate. It is now hoisted out, in both kernels, and fires
identically in all four configurations.

This is the second lesson of Appendix~\ref{app:defects} and the more general one:
a switch meant to isolate a single term has to be checked against what it leaves
alone, not only against what it changes.

\section{The anisotropic filter in code}

the anisotropic filter needed no CUDA change. The rasteriser already accepts a precomputed
covariance, so replacing a scalar variance with a matrix is a change to how the
covariance is built and not to how it is splatted. Three things happen at
construction, once, before the training loop:

\begin{enumerate}[leftmargin=1.4em,itemsep=2pt]
  \item For each primitive, accumulate
        $\Lambda_k = \sum_n w_{nk}(J_nW_n)^{\!\top}\Sigma_{\mathrm{pix}}^{-1}(J_nW_n)$
        over the training cameras that saw it --- the same Jacobian the renderer
        already forms for the EWA transform.
  \item Invert the spectrum to a covariance floor, and floor \emph{that} at
        Mip-Splatting's own scalar value rather than at a re-derivation of it.
        (Taking theirs is what makes Proposition~\ref{prop:reduction} exact in the
        floor as well as in the covariance; re-deriving it differed by $0.5\,\%$
        at the median on the first call, which is enough to seed a different
        densification trajectory.)
  \item Add it to the primitive's own covariance.
\end{enumerate}

$\Lambda_k$ is never touched inside the loop, which is why rendering is
unaffected (\S\ref{sec:fps}) and why the cost is a construction cost
(\S\ref{sec:methodcost}).

\section{Capture protocols}

A protocol is a rule for selecting a subset of the training cameras. The test set
is never touched --- only the capture the model is fitted from. This is what makes
the stress suite cheap: no new images are rendered and no dataset is regenerated.

The subsets are computed in \texttt{tools/camera\_protocols.py} and passed to the
trainer as an index list. \texttt{mixed focal} additionally downsamples the
selected images, which is a different operation and was, for a while,
not being applied at all --- the protocol passed its index list and nothing else,
so it trained a full orbit under another name. \texttt{tools/test\_camera\_subset.py}
now asserts that each protocol selects what it claims to.
